\documentclass{article}
\usepackage{lpalign}
\usepackage{hyperref}

\begin{document}

\paratitle{Non-Linear Program}[NLP]\label{para:NLP}

\begin{namedsubeqs}{NLP}[style=\itshape\roman]
\begin{lpalign}{Minimize}{f(\mathbf{x}) \label{NLP.0}}
  g_i(\mathbf{x}) &\GE 0 &&\forall\ i\in I \label{NLP.1} \\&&
  h_j(\mathbf{x}) &\EQ 0 &&\forall\ j\in J \label{NLP.2} \\&&
  \mathbf{x}      &\GE 0                         \label{NLP.3}
\end{lpalign}
\end{namedsubeqs}

\begin{subequations}\label{std}
\renewcommand{\theequation}{\theparentequation-\roman{equation}}
\begin{align}
  A\mathbf{x} &= \mathbf{b} \label{std:1}\\
  C\mathbf{y} &= \mathbf{d} \label{std:2}
\end{align}
\end{subequations}

Problem \eqref{para:NLP} is a nonlinear program.
Objective \eqref{NLP.0} is followed by constraints (\ref{NLP.1}-\subref{NLP.3}).
Non-negativity is enforced by \subeqref{NLP.3}.

Standard subequations:
\ref{std}, \ref{std:1}, \subref{std:1}, \subeqref{std:2}.

\end{document}
